\section{Extended Cost-Matrix Sensitivity Analysis (Experiments E24 and Task 8)}
\label{app:cost_extension}

To evaluate policy dominance beyond the baseline OCU matrix, Experiment E24 evaluated a 27-point 3-factor grid sweeping attack block reward $r_{\text{block}} \in \{25, 50, 75\}$, quarantine benign penalty $p_{\text{quar}} \in \{7.5, 15.0, 22.5\}$, and restrict benign penalty $p_{\text{rest}} \in \{1.5, 3.0, 4.5\}$ OCU. Across all 27 cost configurations, $P_{2\text{-tier}}$ strictly dominates conservative baseline $P_1$, with net benefit advantage ratios ranging from 1.82x to 2.44x (mean 2.28x).

Furthermore, Task 8 extended the penalty scaling multiplier across $\lambda \in [0.2, 50.0]$ (where Block benign penalty is $50\lambda$). $P_{2\text{-tier}}$ maintains containment dominance over $P_1$ across $\lambda \in [0.2, 30.0]$, with advantage ratios declining from 2.47x (at $\lambda=0.2$) to 1.04x (at $\lambda=30.0$). A critical break-even occurs at $\lambda = 30.8$ (Block penalty = 1,540 OCU): beyond $\lambda > 30.8$ (e.g., at $\lambda = 50.0$, Block penalty = 2,500 OCU, ratio 0.07x), conservative policy $P_1$ becomes strictly optimal because the penalty for collateral blocking of embedded benign flows in Level $L_4$ outweighs containment rewards. Additionally, sweeping the operational regime mixture weight $w_{\text{bal}} \in [0.0, 1.0]$ confirms that $P_{2\text{-tier}}$'s advantage over $P_1$ is bounded between 2.29x (at $w_{\text{bal}} = 1.0$) and 2.48x (at $w_{\text{bal}} = 0.0$).
